\documentclass[10pt,a4paper]{article}
\usepackage[margin=20mm]{geometry}\usepackage[T1]{fontenc}\usepackage{lmodern}
\usepackage{amsmath,booktabs,hyperref}
\setlength{\parindent}{0pt}\setlength{\parskip}{4pt}
\title{2023 Homework 4\\Pre-layout SPICE deliverable}
\author{EE3032 Introduction to VLSI Design}\date{13 October 2026}
\begin{document}\maketitle
\section*{Requested scope}
The original HW4 asks for layout of the HW3 three-bit ripple-carry adder,
extraction, post-layout delay/power and comparison. The user's explicitly
narrower scope stops at the pre-layout \texttt{.sp} circuit. Therefore the
deliverable is the complete HW3 circuit and independently rerun pre-layout
testbench. No layout, DRC, LVS, PEX or post-layout measurement is claimed.
\section*{Circuit and evidence}
\texttt{rca.sp} is byte-equivalent after text line-ending normalization to
the verified HW3 netlist: three 28T full adders, 84 transistors total,
static CMOS with pMOS bodies at VDD and nMOS bodies at ground.
Channel length is $0.18\,\mu$m and the size groups are those in the HW3
report. The four external outputs each drive 1 fF.
\texttt{testbench.sp} uses CIC018 TT, 1.8 V, $25^\circ$C, 20 ps input slew,
2 ps maximum timestep, and the 128-vector, 384 ns HW3 workload.
\texttt{reproduce.py} verifies the netlist hash and launches a separate HSPICE
run. All 128 arithmetic states and output voltage bands pass again.
\begin{center}\begin{tabular}{lr}
\toprule Quantity&Pre-layout result\\\midrule
Verified input states&128/128\\
Worst single-input delay from HW3&833.526988 ps\\
Independently rerun VDD power&177.082000 $\mu$W\\
Independently rerun 384 ns energy&67999.487960 fJ\\\bottomrule
\end{tabular}\end{center}
The delay value comes from HW3's exhaustive 896-transition characterization;
the independent HW4 rerun verifies the functional workload and power.
It does not masquerade as a second exhaustive delay sweep.
\section*{Reference comparison}
The supplied 2023 HW4 reports and pre-layout files of Lin (110501558),
Liao (110501512), Jian (109501570) and Lin (111501559) were inspected.
Their layouts and extracted parasitics belong to their different designs.
The valid full-adder architecture and rail-body connections are useful
references; their post-layout statistics cannot be used as results of this
netlist. A physical comparison would require a new layout and extraction,
outside the requested scope. The compact independent circuit and exhaustive
HW3 proof are retained instead of copying an unrelated PEX file.
\section*{Reproduction and files}
Run \texttt{python reproduce.py --work PATH}; override \texttt{--model PATH}
and \texttt{--hspice PATH} as needed. \texttt{vectors.csv} contains all
four measured output voltages, expected/observed integer and pass flag.
\texttt{validation.json} records the separate run and the shared netlist
SHA256. The local \texttt{D:/cic018.l} model is not redistributed.
\end{document}
